% \iffalse meta-comment
%
% Copyright (C) 2026 by Valentin Boettcher
%
% This file may be distributed and/or modified under the
% conditions of the LaTeX Project Public License, either version 1.3
% of this license or (at your option) any later version.
% The latest version of this license is in:
%    http://www.latex-project.org/lppl.txt
% and version 1.3 or later is part of all distributions of LaTeX
% version 2005/12/01 or later.
%
% \fi
%
% \iffalse
%<*driver>
\documentclass{ltxdoc}
\usepackage{xcolor}
\definecolor{pastelblue}{RGB}{100, 140, 255}
\definecolor{pastelred}{RGB}{255, 100, 100}
\definecolor{pastelgreen}{RGB}{100, 200, 100}
\usepackage{conversationalcomments}
\usepackage{hypdoc}
\EnableCrossrefs
\CodelineIndex
\RecordChanges
\newcommentator{bill}{Bill}{pastelblue}[vale]
\newcommentator{vale}{Vale}{pastelred}[bill]
\newcommentator{alice}{Alice}{pastelgreen}
\begin{document}
  \DocInput{conversationalcomments.dtx}
\end{document}
%</driver>
% \fi
%
% \ProvidesFile{conversationalcomments.dtx}
%             [2026/06/01 v1.0.0 Conversational Review Comments]
%
% \iffalse
%<package>\NeedsTeXFormat{LaTeX2e}
%<package>\ProvidesPackage{conversationalcomments}
%<package>  [2026/06/01 v1.0.0 Conversational Review Comments]
% \fi
%
% \CheckSum{0}
%
% \changes{v1.0.0}{2026/06/01}{Initial documented release}
%
% \GetFileInfo{conversationalcomments.dtx}
%
% \title{The \textsf{conversationalcomments} package\\[1ex]\large \fileversion\ (\filedate)}
% \author{Valentin Boettcher}
%
% \maketitle
%
% \begin{abstract}
%   The \textsf{conversationalcomments} package provides a customizable review mode for \LaTeX\ documents. It is specifically optimized for single- and multi-column layouts (such as RevTeX). Authors can register custom commentators that dynamically generate user macros supporting inline text-linked highlighting, standalone notes, auto-alternating conversation threads, and multi-author conversation blocks.
% \end{abstract}
%
% \section{Introduction}
% When co-authoring scientific papers, leaving clear, context-linked feedback is crucial.
% While standard marginal notes are functional, they often lack conversational structure, leading to scattered replies and confusing discussions.
%
% The \textsf{conversationalcomments} package solves this by establishing a dynamic, conversational review framework.
% Comments are neatly grouped into boxes in the page margins, with colored highlights in the document body and visual connection lines linking each comment directly to its source.
% Multiple authors can participate in structured, alternating discussion threads or collaborative conversation blocks.
%
% \section{Usage Guide}
%
% \subsection{Preamble Setup}
% Simply load the package in your preamble:
% \begin{verbatim}
%   \usepackage{conversationalcomments}
% \end{verbatim}
% By default, this adjusts your document geometry to provide a margin on the right side of the paper for comments.
%
% \DescribeMacro{\setreviewmarginwidth}
% To customize the width of the review margins, use \cs{setreviewmarginwidth}\marg{width}.
% For example, a wider margin accommodates longer comment text:
% \begin{verbatim}
%   \setreviewmarginwidth{1.5in}
% \end{verbatim}
%
% \subsection{Registering Commentators}
%
% \DescribeMacro{\newcommentator}
% The package does not define hardcoded commentators. Instead, you declare your review team in the preamble using:
% \cs{newcommentator}\marg{cmd\_name}\marg{Label}\marg{Color}\oarg{Default Partner}
%
% For example:
% \begin{verbatim}
%   \newcommentator{bill}{Bill}{pastelblue}[vale]
%   \newcommentator{vale}{Vale}{pastelred}[bill]
%   \newcommentator{alice}{Alice}{pastelgreen}
% \end{verbatim}
%
% This generates custom author macros like \cs{bill}, \cs{vale}, and \cs{alice}, as well as their corresponding inline formatters (such as \cs{billreply}, though these are usually invoked automatically).
%
% \subsection{Comment Modes}
%
% \subsubsection{Text-Linked Comments (Highlights)}
% Use the commentator's command with two arguments. The first argument specifies the document text to highlight, and the second is your comment:
% \begin{verbatim}
%   \bill{highlighted text}{This text is blue, and this is my comment.}
% \end{verbatim}
% Here is a live example: \bill{this is a live text-linked comment}{Look! This comment is rendered in the margin with a connecting line back to this text.}
%
% \subsubsection{Standalone Notes}
% To place a floating note in the margin without highlighting any body text, leave the first argument empty:
% \begin{verbatim}
%   \vale{}{This is a floating note next to this line.}
% \end{verbatim}
% \vale{}{This is a live standalone note next to this paragraph.}
%
% \subsubsection{Conversational Threads (Alternating)}
% Provide consecutive arguments to let the package automatically alternate replies between the commentator and their registered default partner:
% \begin{verbatim}
%   \bill{The system dynamics}{Initial comment by Bill}{Vale's reply}{Bill's final follow-up.}
% \end{verbatim}
% Here is a live alternating thread in action: \bill{Live conversation}{Is this standard DTX compilation?}{Yes, we are compiling the package documentation and demonstrating it at the same time!}{Fabulous!}
%
% \subsubsection{Third-Party Nesting}
% A third reviewer can join a two-author thread by nesting their custom command inside a reply argument:
% \begin{verbatim}
%   \bill{Target Text}{Hey Vale, want to grab lunch?}{\vale{Sure!}}{\alice{Can I join too?}}
% \end{verbatim}
% Here is a live example of third-party nesting: \bill{Nesting example}{Lunch today?}{\vale{Yes please!}}{\alice{Count me in!}}
%
% \subsection{Multi-Author Conversations}
%
% \DescribeMacro{\conversation}
% For larger discussions involving more than two people, use the general \cs{conversation}\marg{Target Text}\marg{Body} block. Inside this block, commentator macros automatically switch to behave as inline reply formatters:
% \begin{verbatim}
%   \conversation{Target Text}{
%     \bill{I think we should rewrite this.}
%     \vale{I agree with Bill.}
%     \alice{Let's make sure Valentin reviews it too.}
%   }
% \end{verbatim}
% Here is a live multi-author conversation:
% \conversation{Multi-author demo}{
%   \bill{We are using context-aware macros.}
%   \vale{It works flawlessly inside the DTX driver.}
%   \alice{And the connecting lines look beautiful!}
% }
%
% \StopEventually{\PrintChanges\PrintIndex}
%
% \section{Implementation}
%
% \iffalse
%<*package>
% \fi
%    \begin{macrocode}
\RequirePackage{xcolor}
\RequirePackage{geometry}
\RequirePackage{marginnote}
\RequirePackage{expl3}
\RequirePackage{xparse}
\RequirePackage{tikz}
\usetikzlibrary{calc}

% Review Margin Configuration
\newlength{\reviewmarginwidth}

\newcommand{\setreviewmarginwidth}[1]{
    \setlength{\reviewmarginwidth}{#1}
    \geometry{
        paperwidth=\dimexpr 8.5in + 2\reviewmarginwidth \relax,
        left=\dimexpr 0.75in + \reviewmarginwidth \relax,
        right=\dimexpr 0.75in + \reviewmarginwidth \relax,
        top=1in, bottom=1in,
        marginparwidth=\reviewmarginwidth,
        marginparsep=0.15in
    }
}

% Set default width
\setreviewmarginwidth{2.5in}

\ExplSyntaxOn
\int_new:N \g__comm_count_int
\seq_new:N \l__comm_args_seq
\tl_new:N \l__comm_body_tl
\tl_new:N \l__comm_color_tl
\tl_new:N \l__comm_target_tl
\int_new:N \l__comm_id_int
\bool_new:N \l_comm_first_reply_bool
\bool_new:N \l_comm_in_conversation_bool
\bool_new:N \g_comm_reply_handled_bool

% Dynamic Commentators Registry
\prop_new:N \g_comm_commentators_prop
\tl_new:N \l_initiator_name_tl
\tl_new:N \l_initiator_color_tl
\tl_new:N \l_initiator_partner_id_tl
\tl_new:N \l_partner_name_tl
\tl_new:N \l_partner_color_tl
\tl_new:N \l_curr_name_tl
\tl_new:N \l_curr_color_tl

% Helper to draw the marginal boxes
\cs_new_protected:Nn \__comm_draw_box:nn {
  \group_begin:
  \noindent % Ensure no indent in the margin
  \begin{tikzpicture}[remember~picture, baseline=(commdest-#2.center)]
    \node [
      draw=\l__comm_color_tl,
      fill=white,
      line~width=1.2pt,
      inner~sep=\fboxsep,
      text~width=\dimexpr \marginparwidth - 2\fboxsep - 1.2pt \relax,
      align=left,
      font=\scriptsize\raggedright,
      outer~sep=0pt
    ] (commdest-#2) {
      \bool_set_true:N \l_comm_in_conversation_bool
      \setlength{\parindent}{0pt}%
      #1
    };
  \end{tikzpicture}
  \group_end:
}
\cs_generate_variant:Nn \__comm_draw_box:nn { Vn }

% Render Note (shared positioning and line drawing logic)
\cs_new_protected:Nn \__comm_render_note:n {
  \int_gincr:N \g__comm_count_int
  \int_set_eq:NN \l__comm_id_int \g__comm_count_int
  
  % Marker
  \tikz[remember~picture,overlay,baseline]{
    \node[anchor=base,inner~sep=0pt,outer~sep=0pt](commsrc-\int_use:N \l__comm_id_int){\strut};
  }

  % Detect position and write to aux (for next run)
  \begin{tikzpicture}[remember~picture, overlay]
    \path (commsrc-\int_use:N \l__comm_id_int.center); \pgfgetlastxy{\xn}{\yn}
    \path (current~page.center); \pgfgetlastxy{\xc}{\yc}
    \dim_compare:nNnTF { \xn } < { \xc } {
      \tl_set:Nn \l_tmpa_tl { L }
    }{
      \tl_set:Nn \l_tmpa_tl { R }
    }
    \iow_now:Nx \@auxout {
      \exp_not:N \expandafter \exp_not:N \gdef \exp_not:N \csname commside \int_use:N \l__comm_id_int \exp_not:N \endcsname { \l_tmpa_tl }
    }
  \end{tikzpicture}

  \tl_if_blank:nF { \l__comm_target_tl } { 
    \group_begin:
    \color{\l__comm_color_tl} \l__comm_target_tl 
    \group_end:
  }
  
  % Use the recorded side or default
  \cs_if_exist:cTF { commside \int_use:N \l__comm_id_int } {
    \str_if_eq:eeTF { \use:c { commside \int_use:N \l__comm_id_int } } { L } {
       \reversemarginpar
    }{
       \normalmarginpar
    }
  }{
    \normalmarginpar
  }

  % Placement in margin
  \use:x {
    \exp_not:N \marginnote 
      { \exp_not:N \__comm_draw_box:Vn \exp_not:N \l__comm_body_tl { \int_use:N \l__comm_id_int } } 
      [ #1 ]
  }
  
  \normalmarginpar % Reset

  % Draw the connecting line
  \begin{tikzpicture}[remember~picture,overlay]
    \path (commsrc-\int_use:N \l__comm_id_int.center); \pgfgetlastxy{\commxsrc}{\commysrc}
    \path (commdest-\int_use:N \l__comm_id_int.center); \pgfgetlastxy{\commxdest}{\commydest}
    \ifdim \commxdest > \commxsrc
      \draw[semithick, \l__comm_color_tl, opacity=0.8] 
        ($(current~page.center |- commsrc-\int_use:N \l__comm_id_int.base) + (0.5*\textwidth, 0)$) 
        -- (commdest-\int_use:N \l__comm_id_int.west);
    \else
      \draw[semithick, \l__comm_color_tl, opacity=0.8] 
        ($(current~page.center |- commsrc-\int_use:N \l__comm_id_int.base) - (0.5*\textwidth, 0)$) 
        -- (commdest-\int_use:N \l__comm_id_int.east);
    \fi
  \end{tikzpicture}
}

% Core alternating processing logic
\cs_generate_variant:Nn \prop_get:NnN { NoN }

% Dynamic evaluation macro for a reply
\box_new:N \l_comm_scratch_box

\cs_new_protected:Nn \__comm_eval_reply:nn {
  \bool_gset_false:N \g_comm_reply_handled_bool
  \group_begin:
    % Dry run in a scratch vbox to detect commentator macros
    \vbox_set:Nn \l_comm_scratch_box {
      \hsize=\dimexpr \marginparwidth - 2\fboxsep - 1.2pt \relax
      #2
    }
  \group_end:
  
  \bool_if:NTF \g_comm_reply_handled_bool {
    % It was self-formatted: execute it directly in the main stream
    #2
  }{
    % Plain text: format with default author
    \use:c { #1 reply } { #2 }
  }
}

\cs_new_protected:Nn \__comm_process:nn {
  \bool_gset_true:N \l_comm_first_reply_bool
  
  \seq_get_left:NN \l__comm_args_seq \l_tmpa_tl
  \seq_pop_left:NN \l__comm_args_seq \l_tmpb_tl
  
  \tl_set_eq:NN \l__comm_target_tl \l_tmpa_tl

  \seq_if_empty:NTF \l__comm_args_seq {
    \tl_set_eq:NN \l__comm_body_tl \l_tmpa_tl
    \tl_clear:N \l__comm_target_tl
  }{
    \tl_clear:N \l__comm_body_tl
    \prop_get:NnN \g_comm_commentators_prop { #1/partner } \l_initiator_partner_id_tl
    
    \int_step_inline:nn { \seq_count:N \l__comm_args_seq } {
      \int_if_odd:nTF { ##1 } { 
        \tl_set:Nn \l_curr_id_tl { #1 }
      }{ 
        \tl_set_eq:NN \l_curr_id_tl \l_initiator_partner_id_tl
        \tl_if_empty:NT \l_curr_id_tl {
          \tl_set:Nn \l_curr_id_tl { #1 }
        }
      }
      
      \tl_put_right:Nx \l__comm_body_tl { 
        \exp_not:N \__comm_eval_reply:nn { \l_curr_id_tl } { \seq_item:Nn \l__comm_args_seq { ##1 } } 
      }
    }
  }

  \prop_get:NnN \g_comm_commentators_prop { #1/color } \l__comm_color_tl
  \__comm_render_note:n { #2 }
}

% Dynamic Commentator Registration
\cs_new_protected:Nn \__comm_register_commentator:nnnn {
  \prop_gput:Nnn \g_comm_commentators_prop { #1/name } { #2 }
  \prop_gput:Nnn \g_comm_commentators_prop { #1/color } { #3 }
  \prop_gput:Nnn \g_comm_commentators_prop { #1/partner } { #4 }

  % Main alternating command: \<id> (e.g. \bill)
  \exp_args:Nc \NewDocumentCommand { #1 } { m +g +g +g +g +g +g +g O{0pt} } {%
    \bool_if:NTF \l_comm_in_conversation_bool {%
      % Inside a conversation: act as inline reply (format first argument)
      \use:c { #1 reply } { ##1 }%
    }{%
      % Thread starter
      \seq_clear:N \l__comm_args_seq%
      \seq_put_right:Nn \l__comm_args_seq { ##1 }%
      \IfValueT { ##2 } { \seq_put_right:Nn \l__comm_args_seq { ##2 } }%
      \IfValueT { ##3 } { \seq_put_right:Nn \l__comm_args_seq { ##3 } }%
      \IfValueT { ##4 } { \seq_put_right:Nn \l__comm_args_seq { ##4 } }%
      \IfValueT { ##5 } { \seq_put_right:Nn \l__comm_args_seq { ##5 } }%
      \IfValueT { ##6 } { \seq_put_right:Nn \l__comm_args_seq { ##6 } }%
      \IfValueT { ##7 } { \seq_put_right:Nn \l__comm_args_seq { ##7 } }%
      \IfValueT { ##8 } { \seq_put_right:Nn \l__comm_args_seq { ##8 } }%
      \__comm_process:nn { #1 } { ##9 }%
    }%
  }

  % Inline formatter: \<id>reply (e.g. \billreply)
  \exp_args:Nc \NewDocumentCommand { #1 reply } { m } {%
    \bool_gset_true:N \g_comm_reply_handled_bool%
    \bool_if:NTF \l_comm_first_reply_bool {%
      \bool_gset_false:N \l_comm_first_reply_bool%
    }{%
      \par\vspace{3pt}%
    }%
    \textcolor { #3 } { \textbf { #2 : } ~ ##1 }%
  }
}

% Public registration API
\NewDocumentCommand{\newcommentator}{m m m O{}}{
  \__comm_register_commentator:nnnn { #1 } { #2 } { #3 } { #4 }
}

% Multi-author explicit conversation command
\NewDocumentCommand{\conversation}{m +m O{0pt}}{
  \group_begin:
  \bool_gset_true:N \l_comm_first_reply_bool
  \tl_set:Nn \l__comm_target_tl { #1 }
  \tl_set:Nn \l__comm_body_tl { #2 }
  \tl_trim_spaces:N \l__comm_body_tl
  \tl_set:Nn \l__comm_color_tl { gray } % Neutral rim for multi-author
  \__comm_render_note:n { #3 }
  \group_end:
}

\ExplSyntaxOff
%    \end{macrocode}
%
% \iffalse
%</package>
% \fi
%
% \Finale
\endinput
